\section{Kapitel 6}
\subsection{Teilaufgabe 1}
\subsubsection{Aufgabenstellung}
Im Grunde ist es die Selbe Aufgabe wie aus Kapitel 4, Teilaufgabe 2. Doch jetzt solle es auch
für nxn Matrizen funktionieren. Die grö\ss\space gibt an ende der Nutzer ein. Zusätzlich soll
noch die Multiplikation der Matrizen auch mit while und do-while gelöst werden.

\subsubsection{Anforderungsdefinition}
\begin{enumerate}
	\item Unser Input für die Grö\ss\space der nxn Matrix.
	\item Multiplikation mit for, while, und do-while.
\end{enumerate}

\subsubsection{Entwurf}
Zuerst wird der Nutzer aufgefordert eine Zahl einzugeben. Anschließend wird in einer do while 
Schleife geprüft ob die Nutzereingabe valide ist. Es wird einmal geprüft ob es sich um Buchstaben handelt
und ob die Zahl kleiner als 2 ist. Wenn der Fall nicht eintritt, wird die Schleife beendet und es werden zwei
nxn Matrizen erzeugt. Der Rest ist identisch wie aus Kapitel 4, nur das hier noch für die Berechnung 
eine weitere while und do while Schleife verwendet wird
\begin{center}
	\includegraphics[width=0.99\textwidth]{uml/uml_c6_p1.pdf}
\end{center}

\subsubsection{Quelltext}
\paragraph{Matrizen.java}\
\lstinputlisting[language = Java , frame = trBL , escapeinside={(*@}{@*)}]{../chapter_06/src/main/java/chapter_06/Matrizen.java}

\subsubsection{Testdokumentation}
Bei Zahlen die kleiner als 2 waren, sowie Buchstaben, kam es zu einer entsprechenden Fehlermeldung. Anschlie\ss end konnte man den Wert erneut eintippen.
Bei den Ergebnissen der Matrizen wurden die Werte mit einem Onlinerechner verglichen, hierbei hat
alles übereingestimmt.

\subsubsection{Benutzungshinweise}
Nach dem aufrufen des Programms, wird der Nutzer aufgefordert eine Zahl einzugeben.
Diese muss grö\ss er als ein sein.

\subsubsection{Anwendungsbeispiel}
Nach dem man das Programm gestartet hat, sollte folgende Ausgabe erscheinen:
\begin{lstlisting}[frame = trBL , escapeinside={(*@}{@*)}]
[sebastian@laptop bin]$ java Matrizen 
Dieses Programm berechnet eine zufällig erstellte nxn Matrix
Geben sie n an: -45
Die Zahl muss größer gleich 2 sein
test
Es dürfen nur Zahlen verwendet werden
5
Matrix A:
78	88	96	91	50	
17	16	14	17	77	
34	1	45	21	42	
63	1	92	76	57	
84	59	13	85	46	

Matrix B:
38	36	45	22	18	
96	4	0	93	79	
86	38	92	50	92	
1	54	63	71	10	
4	62	91	31	55	

Addition von A und B:
116	124	141	113	68	
113	20	14	110	156	
120	39	137	71	134	
64	55	155	147	67	
88	121	104	116	101	

Multiplikation von A und B:
For Schleife
19959	14822	22625	22711	20848	
3711	6900	10131	6156	7263	
5447	6676	10815	5884	7351	
10706	13406	21274	13242	13572	
10243	11196	14517	15446	10749	

While Schleife
19959	14822	22625	22711	20848	
3711	6900	10131	6156	7263	
5447	6676	10815	5884	7351	
10706	13406	21274	13242	13572	
10243	11196	14517	15446	10749	

Do-While Schleife
19959	14822	22625	22711	20848	
3711	6900	10131	6156	7263	
5447	6676	10815	5884	7351	
10706	13406	21274	13242	13572	
10243	11196	14517	15446	10749
[sebastian@laptop bin]$ 
\end{lstlisting}